%==============================================================================
\section{Implementation --- Triển khai học máy trong thực tế}
%==============================================================================

\begin{dongco}[title={target}]
Bài này mô tả các bước triển khai học máy (implementation) trong thực tế: từ dữ liệu, trực quan hoá, feature engineering đến chọn mô hình, đánh giá, tuning và triển khai/giám sát.
\end{dongco}

\subsection{Các bước triển khai học máy}

\subsubsection{Data Collection and Preparation}
\begin{phuongphap}[title={Thu thập và chuẩn bị dữ liệu}]
Bước đầu tiên là thu thập dữ liệu để huấn luyện và kiểm thử mô hình.
Dữ liệu cần \textbf{liên quan} bài toán.
Sau khi thu thập, cần tiền xử lý và làm sạch để loại bỏ dữ liệu bị thiếu, nhiễu, trùng lặp, \ldots
\end{phuongphap}

\subsubsection{Data Exploration and Visualization}
\begin{phuongphap}[title={Khám phá và trực quan hoá dữ liệu}]
Bước tiếp theo là khám phá và trực quan hoá để hiểu cấu trúc và nhận diện pattern/trend.
Các công cụ như Matplotlib và Seaborn để vẽ histogram, scatter plots, heatmaps.
\end{phuongphap}

\subsubsection{Feature Selection and Engineering}
\begin{phuongphap}[title={Chọn feature và tạo feature}]
Chọn các feature liên quan bài toán hoặc tạo feature mới từ feature sẵn có để tăng độ chính xác mô hình.
\end{phuongphap}

\subsubsection{Model Selection and Training}
\begin{phuongphap}[title={Chọn mô hình và huấn luyện}]
Sau khi dữ liệu chuẩn bị và feature được chọn/tạo, bước tiếp theo là chọn thuật toán phù hợp để huấn luyện.
Chia train/test và fit mô hình trên train.
Có thể dùng nhiều thuật toán như linear regression, logistic regression, decision trees, random forests, SVM, neural networks, \ldots
\end{phuongphap}

\subsubsection{Model Evaluation}
\begin{phuongphap}[title={Đánh giá mô hình}]
Sau khi huấn luyện, cần đánh giá hiệu năng.
Metric như accuracy, precision, recall, F1-score; và có thể dùng cross-validation để kiểm thử khả năng tổng quát.
\end{phuongphap}

\subsubsection{Model Tuning}
\begin{phuongphap}[title={Điều chỉnh siêu tham số}]
Có thể cải thiện hiệu năng bằng điều chỉnh siêu tham số (các tham số không học từ dữ liệu mà do người dùng đặt).
Tìm bộ tối ưu bằng grid search và random search.
\end{phuongphap}

\subsubsection{Deployment and Monitoring}
\begin{phuongphap}[title={Triển khai và giám sát}]
Sau khi huấn luyện và tuning, triển khai mô hình vào production.
Triển khai gồm tích hợp mô hình vào quy trình/hệ thống.
Sau đó cần giám sát thường xuyên để đảm bảo mô hình tiếp tục hoạt động tốt và phát hiện vấn đề cần xử lý.
\end{phuongphap}

\begin{luuy}[title={Kết luận}]
Mỗi bước cần công cụ/kỹ thuật khác nhau.
Triển khai thành công cần kết hợp kỹ năng kỹ thuật và hiểu bối cảnh kinh doanh.
\end{luuy}

\subsection{Chọn ngôn ngữ và IDE để phát triển ML}

\begin{dongco}[title={Chọn công cụ theo năng lực và bối cảnh}]
Để phát triển ứng dụng ML, bạn cần quyết định nền tảng, IDE và ngôn ngữ.
Nhiều lựa chọn có thể đáp ứng nhu cầu vì đều có hiện thực các thuật toán AI/ML phổ biến.
\end{dongco}

\begin{phuongphap}[title={Nếu tự phát triển thuật toán}]
Gợi ý cân nhắc:
\begin{itemize}
  \item Ngôn ngữ bạn thành thạo.
  \item IDE bạn quen dùng.
  \item Nền tảng triển khai (nhiều nền tảng miễn phí; một số có phí license theo mức sử dụng).
\end{itemize}
\end{phuongphap}

\subsubsection{Language Choice}
\begin{phuongphap}[title={Ngôn ngữ hỗ trợ ML}]
Python, R, Matlab, Octave, Julia, C++, C\ldots
\end{phuongphap}

\subsubsection{IDEs}
\begin{phuongphap}[title={IDE hỗ trợ ML}]
R Studio, PyCharm, iPython/Jupyter Notebook, Julia, Spyder, Anaconda, Rodeo, Google Colab\ldots
Lưu ý: mỗi IDE có ưu/nhược, nên thử trước khi chọn một IDE chính.
\end{phuongphap}
